\chapter{Conclusion and Future Work}
\label{ch:conclusion_and_future_work}
\label{sec:conclusion_and_future_work}

\section{Conclusion}
\label{sec:conclusion}

This thesis investigates \gls{id}-based synchronization for edge-based \glspl{dt} in bandwidth-constrained environments. Beyond reducing traffic volume, the \gls{id} mechanism can reduce end-to-end latency by replacing full payload transfer with \enquote{state verification}. The results support the following conclusions:

First, the \gls{id} protocol reduces the sensitivity of synchronization latency to state dimension in bandwidth-limited regimes. Experiments show that for $B < 1$ Mbps, \gls{id} yields lower end-to-end latency than full-state transmission. For high-dimensional sensor data such as \gls{lidar}, the additional encoding overhead is typically outweighed by reduced fallback transmission cost, improving freshness by avoiding unnecessary payload transfers.

Second, for the evaluated software stack, truncated hashing based on \glsentryshort{sha}-256 is typically a more practical engineering choice than algebraic coding via \gls{rsid} for edge deployments. The main driver is implementation efficiency: optimized \glsentryshort{sha}-256 backends achieve low and largely payload-independent encoding latency, whereas \gls{rsid} incurs higher computation that depends on the symbolization of the payload.

\gls{rsid} nevertheless remains relevant in two regimes. The first is high-integrity applications where rigorous mathematical error bounds are mandatory (i.e., a provable upper bound on failure probability under the stated modeling assumptions). The second is transmission-congested, fallback-heavy operation, where full-state recovery is triggered frequently and the resulting transmission delay dominates the end-to-end latency; in this regime, marginal compute-time differences between encoding mechanisms become comparatively less influential.

Third, \gls{kid} implements the shift from \enquote{bitwise exactness} to \enquote{semantic compliance}. The Pareto analysis shows that introducing a semantic tolerance radius $R$ allows the system to trade reduced verification frequency (and thus reduced fallback traffic) against increased acceptance-set size. In the evaluated sweeps, a moderate tolerance improves speedup substantially, while configurations with very loose tolerance (e.g., $R=30$) do not appear among the non-dominated points due to a saturation in speedup relative to the growth of the acceptance set.

As a practical guideline for the studied noise levels, $R\approx 20$ behaves as a knee point: further increases yield limited marginal speedup while reducing verifier selectivity. Here, selectivity refers to the verifier's physical sensitivity in distinguishing negligible drift (treated as semantically admissible) from critical deviations that should trigger recovery; as the acceptance set expands, a larger fraction of off-nominal states becomes indistinguishable from admissible noise at the verifier interface.

Fourth, the dynamic evaluation highlights time-correlated risk in the presence of probabilistic verification. Under random-walk drift, weaker verifiers can turn isolated collisions into persistent, undetected boundary breaches once $|E_t|>R$. The throughput analysis based on the \gls{sri} illustrates a two-tier control logic: the physical tolerance radius $R$ regulates the base synchronization rhythm (traffic sparsification), while verifier entropy $n_{ver}$ governs recovery reliability by suppressing collision-induced misses.

This also exposes an entropy-saturation feasibility constraint. Let $K\triangleq |\mathcal{A}_R|$ denote the acceptance-set size induced by the physical radius $R$ at quantization step $\Delta$ (e.g., in 1D, $K=2\lfloor R/\Delta\rfloor+1$). Under the accept-set approximation, the miss probability scales as $p_{miss}\approx K/2^{n_{ver}}$. If $K$ is comparable to or exceeds $2^{n_{ver}}$ (e.g., large $R$ at fine $\Delta$ with small $n_{ver}$), the verifier space becomes saturated and reliable detection cannot be guaranteed. Finally, the dynamic metrics show an inherent trade-off: increasing $R$ can reduce the typical duration of undetected intervals but increases the maximum physical drift accumulated during such events.

Overall, the evaluation quantifies the coupled trade-offs among computation, communication, and control fidelity, and yields design guidance for semantic synchronization policies in which tolerance radius $R$, acceptance-set size $K$, verifier entropy ($n_{ver}$), and quantization ($\Delta$) are selected jointly under bandwidth and safety constraints.

\section{Future Work}
\label{sec:future_work}

While this work has characterized the latency and reliability of \gls{kid}, several avenues remain to extend its theoretical depth and practical applicability in next-generation networks.

\subsection{Energy-Aware Identification Policies}
Current evaluations utilize traffic volume and latency as proxies for system efficiency. However, for battery-powered \gls{iot} devices, energy consumption is often the dominant constraint. Future work should explicitly model the energy--delay product, balancing the computational energy cost of hashing against the transmission energy cost of radio interfaces. By integrating the First-Order Radio Model~\cite{heinzelman2000energy} with real-time power profiling of edge hardware, we aim to derive an energy-optimal policy that dynamically adapts the tolerance radius $R$ (and thus the induced acceptance-set size $K$) and verifier length $n_{ver}$ based on the device's remaining battery life.

\subsection[Integration with 6G Physical Layer Security]{Integration with \glsentryshort{sixg} Physical Layer Security}
The current implementation relies on synchronized \glspl{prng}, which introduces a dependency on upper-layer clock synchronization protocols like \gls{ptp}. To align with the architecture of \gls{sixg} networks, future iterations will investigate the integration of \gls{plkg}. By exploiting the channel reciprocity in \gls{tdd} systems, the robot and the \gls{dt} can extract common randomness directly from the \gls{csi}. This approach would turn channel noise into a synchronization resource, eliminating the need for clock synchronization and enhancing robustness against timing attacks in harsh industrial environments.

\subsection{Scalability to Multi-Agent Systems}
This thesis focused on a single robot--\gls{dt} loop. However, industrial scenarios often involve swarms of robots sharing a physical workspace. A critical extension is the development of semantic multiple access schemes, where multiple robots superimpose their identification codes on a shared channel. Future work will explore algebraic constructions that allow the \gls{dt} to verify the aggregate safety of the swarm directly in the coded domain, checking for \enquote{Semantic Inter-Robot Collisions} without decoding individual states, thereby scaling the \gls{kid} paradigm to massive multi-agent systems.

\subsection[Predictive DTs and Adaptive K-ID]{Predictive \glsentryshortpl{dt} and Adaptive \glsentryshort{kid}}
The dynamic drift evaluation in this thesis uses a deliberately simple predictor (\gls{zoh}) to isolate protocol-induced effects. A natural next step is to integrate stronger predictors such as Kalman filtering or learned kinematic models, and to study how prediction quality reshapes both the masking risk (through the breach rate) and the collision risk (through the verification rate). In addition, future work should explore adaptive policies that tune $(R,n_{ver})$ online based on observed breach statistics and safety budgets, tightening the tolerance under persistent drift and relaxing it during stable operation to reduce bandwidth.

\subsection{Decoder Compute, Scheduling, and Quantization Sensitivity}
Beyond bandwidth and latency, practical deployments must also account for \glsentryshort{dt}-side compute and memory. Since the acceptance-set size scales as $K=|\mathcal{A}_R|\approx 2\lfloor R/\Delta\rfloor+1$ in 1D, a direct implementation precomputes $\mathcal{V}_R$ by hashing $\mathcal{O}(K)$ states per \glsentryshort{dt} update, and stores up to $|\mathcal{V}_R|\le K$ verifier codewords. Future work should therefore characterize end-to-end performance under explicit decoder compute constraints (e.g., bounded hash evaluations per tick), and study algorithmic accelerations (incremental updates, caching, \gls{gpu} offload).

In addition, the dynamic setting naturally suggests a scheduling-centric view: rather than only reporting per-tick risk, one can characterize the distribution of full-state update intervals (e.g., via \glspl{sri}, survival functions, and \gls{cdf} metrics) to answer questions like ``with 99\% probability, how long does the channel remain idle?'' Finally, because $\Delta$ is tied to sensor precision and \gls{dt} compute budget, a systematic sensitivity analysis over $\Delta$ is needed to guide robust configuration across heterogeneous sensing modalities.